% !TeX program = xelatex
%%%%%%%%%%%%%%%%%%%%%%%%%%%%%%%%%%%%%%%%%
% Medium Length Professional CV
% Cleaned Version
%%%%%%%%%%%%%%%%%%%%%%%%%%%%%%%%%%%%%%%%%

\documentclass{resume}

%----------------------------------------------------------------------------------------
% PACKAGES AND DOCUMENT CONFIGURATION
%----------------------------------------------------------------------------------------

\usepackage[left=0.4in,top=0.3in,right=0.4in,bottom=0.3in]{geometry}
\usepackage{ctex}
\setmainfont{LinLibertine_R.otf}[
    BoldFont=LinLibertine_RB.otf,
    ItalicFont=LinLibertine_RI.otf,
    BoldItalicFont=LinLibertine_RBI.otf
]
\setCJKmainfont{Songti SC}
\usepackage{enumitem}
\usepackage{hyperref}

\hypersetup{
    colorlinks=false,
    pdfborder={0 0 0}
}

\setlist[itemize]{
    leftmargin=1.5em,
    itemsep=1pt,
    topsep=2pt,
    parsep=0pt,
    partopsep=0pt
}

\newcommand{\subsectiontitle}[1]{%
    \vspace{2pt}
    \textbf{#1}
    \vspace{2pt}
}

\newenvironment{publicationentry}
{\begin{list}{}{
    \setlength{\leftmargin}{1.2em}
    \setlength{\rightmargin}{0em}
    \setlength{\topsep}{2pt}
    \setlength{\partopsep}{0pt}
    \setlength{\parsep}{0pt}
    \setlength{\itemsep}{0pt}
}\item[]}
{\end{list}}

\name{倪培洹}

\address{北京市海淀区中关村东路95号，中国科学院自动化研究所}

\address{
+86 158 0375 7631 \\
\href{mailto:nipeihuan24@mails.ucas.ac.cn}
{nipeihuan24@mails.ucas.ac.cn} \\
\href{https://peihuanni.github.io}
{peihuanni.github.io}
}

\begin{document}

\begin{rSection}{研究方向}
面向大模型训练与推理的人工智能基础设施和分布式系统；大语言模型与视觉--语言--动作模型的高效推理架构；硬件加速器与具身人工智能。
\end{rSection}

\begin{rSection}{博士申请意向}
正在寻找2027年入学的国内外博士机会，研究方向包括机器学习系统（训练与推理加速）、人工智能基础设施和计算机体系结构。
\end{rSection}

%----------------------------------------------------------------------------------------
% EDUCATION
%----------------------------------------------------------------------------------------

\begin{rSection}{教育经历}

\textbf{吉林大学，学士}
\hfill
2020年9月 -- 2024年6月

电子科学与工程学院，微电子科学与工程
\hfill
导师：阎凯歌

CET-4：603，CET-6：553

\vspace{3pt}

\textbf{中国科学院大学，硕士研究生}
\hfill
2024年9月 -- 2027年6月

人工智能学院，人工智能
\hfill
导师：程健、李钢

中国科学院自动化研究所，智能计算与学习实验室（CLab）

GPA：3.81

% \vspace{3pt}

% \textbf{南京大学，博士研究生（预计）}
% \hfill
% 2027年9月 -- 2031年6月

% 计算机学院，计算机科学与技术
% \hfill
% 导师：田臣、王智彬

\end{rSection}

%----------------------------------------------------------------------------------------
% RESEARCH EXPERIENCE
%----------------------------------------------------------------------------------------

\begin{rSection}{科研经历}

\textbf{第一作者论文}

\begin{publicationentry}
\textbf{Peihuan Ni}, Zitao Mo, Tielong Liu, Hongli Wen, Zeyu Zhu,
Minnan Pei, Junwen Si, Weifan Guan, Peisong Wang, Qinghao Hu,
Gang Li and Jian Cheng.
\textbf{APEX: Integer-only Non-linear Function Approximation for Efficient
Cross-Modal Inference.}
In \textbf{DATE 2026}.
\textbf{(CCF-B，计算机体系结构顶级会议)}

\begin{itemize}
    \item 发现不同模态模型在量化后的数据分布存在显著差异，使其对非线性函数近似误差表现出不同敏感性，而统一保留高精度会导致计算位宽和硬件开销过高。
    \item 提出统一的整数非线性计算架构，融合多种常用非线性算子的计算图，并通过比特级剪枝减少计算位宽，同时保持模型精度。
    \item 在语言模型和视觉模型上分别提升最高约0.7\%和1.3\%的任务精度，并在多模态模型上实现近乎无损推理。
    \item 相较已有方法，面积降低1.73--8.71$\times$，功耗降低1.21--10.83$\times$。
\end{itemize}
\end{publicationentry}

\begin{publicationentry}
\textbf{Peihuan Ni}, Zeyu Zhu, Weifan Guan, Gang Li, Zitao Mo, Tielong Liu,
Junwen Si, Qinghao Hu, Peisong Wang and Jian Cheng.
\textbf{ACE-VLA: Action-Centric Evaluation for Token Pruning in
Vision-Language-Action Models.}
In \textbf{CoRL 2026}.
\textbf{(智源A类会议，具身智能顶级会议)}

\begin{itemize}
    \item 指出现有VLA剪枝方法主要沿用VLM中的文本--视觉注意力评价方式，未充分利用动作输出中包含的任务阶段和运动语义信息。
    \item 提出以动作输出为中心的梯度评价指标，并使用基于运动轨迹的动态Token池维护策略降低重要性评估开销。
    \item 从运动学角度建模机器人的运动趋势，而非仅依据瞬时运动状态，实现动态剪枝比例调整。
    \item 在$\pi_{0.5}$和$\pi_0$模型上分别提升约1.4\%和2.2\%的任务成功率，推理速度提高$1.33$--$1.57\times$，并在OpenVLA-OFT上取得优于现有基线的综合性能。
    \item 通过真机实验验证方法在真实场景中的可迁移性，并提升VLA剪枝过程的可解释性。
\end{itemize}
\end{publicationentry}

\vspace{2pt}

\textbf{共同作者论文}

\begin{publicationentry}
Zeyu Zhu, Gang Li, Minnan Pei, Zitao Mo, \textbf{Peihuan Ni},
Peisong Wang, Tielong Liu and Jian Cheng.
\textbf{KL-MoE: A Hierarchical MoE Pruning Framework Exploiting KL Divergence.}
In \textbf{DAC 2026}.
\textbf{(CCF-A，并行与分布计算顶级会议)}
\end{publicationentry}

\begin{publicationentry}
Tielong Liu, Zeyu Zhu, Gang Li, Zitao Mo, Chen BAI, Junwen Si,
\textbf{Peihuan Ni}, Minnan Pei, Peisong Wang, Zhuoran Song,
Xiaoyao Liang and Jian Cheng.
\textbf{FLUTE: A LUT-Based Floating-Point Engine for Efficient LLM
Acceleration Exploiting Core-Tail Hierarchical Alignment.}
In \textbf{MICRO 2026}.
\textbf{(CCF-A，计算机体系结构顶级会议)}
\end{publicationentry}

\begin{publicationentry}
Junwen Si, Gang Li, Tielong Liu, \textbf{Peihuan Ni},
Minnan Pei, Zeyu Zhu, Zitao Mo, Peisong Wang,
Zhongfeng Wang and Jian Cheng.
\textbf{DOMI: A Dataflow Optimization Framework for Efficient Mamba
Inference on FPGAs.}
Submitted to \textbf{TCAD}.
\textbf{(CCF-A，计算机体系结构顶级期刊)}
\end{publicationentry}

\begin{publicationentry}
Zeyu Zhu, Gang Li, Peisong Wang, \textbf{Peihuan Ni},
Tielong Liu, Minnan Pei, Junwen Si, Zhuoran Song,
Xiaoyao Liang and Jian Cheng.
\textbf{DALI: A Workload-Aware Offloading Framework for Efficient MoE
Inference on Local PCs.}
Submitted to \textbf{ATC 2026}.
\textbf{(CCF-A，并行与分布计算顶级会议)}
\end{publicationentry}

\end{rSection}

%----------------------------------------------------------------------------------------
% INTERNSHIP
%----------------------------------------------------------------------------------------

\begin{rSection}{实习经历}

\textbf{华为，ICT BG}
\hfill
\textbf{Ascend 950 NPU大模型推理系统加速}
\hfill
\textbf{2026年7月 -- 2027年7月}

\begin{itemize}
    \item 面向Ascend 950 NPU开展大模型推理系统分析与性能优化，重点研究算子调度、片上数据流和计算资源利用率。
    \item 分析大模型推理过程中Cube、Vector及存储层次之间的性能瓶颈，并探索适用于MoE和长序列推理的系统级优化方法。
\end{itemize}

\end{rSection}

%----------------------------------------------------------------------------------------
% PROJECTS
%----------------------------------------------------------------------------------------

\begin{rSection}{项目经历}

\textbf{VLA模型量化压缩与FPGA端侧部署}
\hfill
\textbf{项目负责人}

\begin{itemize}
    \item 设计W4A8量化方案，其中权重采用4-bit量化，激活值和KV Cache采用8-bit量化，并结合纯整数非线性近似实现近乎无损的端侧推理。
    \item 实现基于FP16--INT8 MAC的量化与反量化融合算子，减少中间数据回写和片外访存，提高激活值的片上驻留比例。
    \item 复用FP16--INT8 MAC完成RoPE等算子的计算，并调整部分计算顺序，提高计算单元和片上数据的复用效率。
    \item 面向W4A8混合精度矩阵乘法设计比特串行计算单元，以及支持不同注意力头维度的KV重排序模块和BRAM到计算阵列的Crossbar数据映射。
    \item 设计统一的纯整数非线性计算单元，在降低硬件面积和功耗的同时支持多类非线性算子，相关成果发表于DATE 2026。
    \item 参与VLIW编译器设计，负责片上数据流、指令处理器和算子控制模块；通过地址偏移和多种偏移模式，将单次模型推理所需指令数降低32$\times$以上。
    \item 在GPU上实现Bit-Accurate仿真器，并完成与Vivado模块级仿真及FPGA整网运行结果的逐级对齐。
    \item 完成DMA访存通路调试、HBM带宽优化、VHK158开发板数据通路调试、时序收敛及PetaLinux系统制作，实现模型在FPGA上的完整部署。
\end{itemize}

\vspace{3pt}

\textbf{基于Ascend A5 NPU的大模型分析与推理优化}
\hfill
\textbf{项目负责人}

\begin{itemize}
    \item 分析Ascend架构中矩阵计算单元与向量计算单元算力不对称导致的性能瓶颈，定位MoE推理中的非线性算子、数据搬运和计算资源空闲问题。
    \item 探索矩阵乘法与非线性计算的细粒度流水和重叠执行方案，以改善Cube与Vector之间的负载均衡并提升整体硬件利用率。
\end{itemize}

\end{rSection}

%----------------------------------------------------------------------------------------
% AWARDS
%----------------------------------------------------------------------------------------

\begin{rSection}{奖项荣誉}

\textbf{综合荣誉}

中国科学院大学三好学生（2025）

\vspace{3pt}

\textbf{竞赛获奖}

英特尔杯全国大学生电子设计竞赛邀请赛三等奖（2022）

全国大学生数学建模竞赛二等奖（2022）

中国大学生数学竞赛三等奖（2023）

\end{rSection}

%----------------------------------------------------------------------------------------
% ACADEMIC SERVICE
%----------------------------------------------------------------------------------------

\begin{rSection}{社会服务}

WACV、EMNLP、CoRL、HPCA审稿人

\end{rSection}

\end{document}